
Among 115 Papers With Code benchmarks, a benchmark's competitive landscape undergoes a discrete phase transition --- not a gradual fade --- at approximately 39 competing papers. Before this threshold, performance scores vary nearly $4\times$ more than the global baseline as the community searches for winning approaches. After it, variance collapses by 80\% relative to the pre-breakpoint segment. This structural break, detectable in two minutes of computation, offers a principled, quantitative answer to a question the ML evaluation community has long debated informally: when should we retire a benchmark?

The need for benchmark retirement criteria is not merely academic. When a benchmark saturates --- when top models differ by fractions of a percent --- it can no longer rank research contributions meaningfully. Continuing to optimize such benchmarks exemplifies Goodhart's Law at scale: once a measure becomes a target, it ceases to be a good measure [Goodhart1975]. Yet retirement decisions in the ML community remain largely informal, driven by qualitative consensus rather than quantitative evidence [Bowman2021].

Prior work has characterized benchmark saturation in aggregate. Liao et al. [Liao2022] documented near-saturation trends across 3,765 benchmarks using time-based metrics and coefficient of variation. The S\_index [Akhtar2026] provides a composite saturation score for LLM benchmarks. These contributions establish that saturation is a systemic phenomenon --- but they do not answer the retirement question quantitatively: at what point in a benchmark's publication history does saturation onset occur?

This work addresses that gap by reframing benchmark saturation as a structural break detection problem. The key insight is that benchmark performance variance does not decline gradually --- it undergoes a statistically significant regime shift at a detectable paper\_count threshold. Rather than fitting a monotonic trend to CoV-vs-paper\_count data, PELT change-point detection [Killick2012] is applied to OLS-detrended residual CoV, directly targeting the transition between exploration and saturation regimes. Two independent tests --- a permutation test on breakpoint position (p=0.035) and a piecewise F-test comparing one-segment vs. two-segment linear models (p=0.0021) --- both confirm this transition at paper\_count\* $\approx$ 39.

\textbf{Contributions.}

\begin{enumerate}
\item \textbf{Empirical (novel threshold).} The first detected structural break in PwC benchmark residual CoV at paper\_count\* $\approx$ 39 (N=115 benchmarks, Aug 2026 snapshot), providing a quantitative, paper-count-based candidate threshold for benchmark retirement flagging. This threshold is not available in prior work.
\end{enumerate}

\begin{enumerate}
\item \textbf{Empirical (two-regime structure).} Characterization of two regimes: a high-variance exploration phase (pre-breakpoint variance 3.$81\times$ global baseline, F p=0.0009) and a low-variance saturation phase (post-breakpoint variance 0.$20\times$ pre-breakpoint, Brown-Forsythe p=0.0099). Three experiments with distinct test designs and null hypotheses converge on the same two-regime narrative.
\end{enumerate}

\begin{enumerate}
\item \textbf{Methodological.} Demonstration that PELT on OLS-detrended residual CoV is a viable, computationally efficient, and interpretable approach for detecting saturation onset in benchmark leaderboard data. The methodology is dataset-agnostic, fully automated, and reproducible.
\end{enumerate}

The paper is organized as follows. Section 2 reviews related work. Section 3 describes the methodology. Section 4 presents experimental setup. Section 5 reports results. Section 6 discusses interpretation and limitations. Section 7 concludes.

